\chapter{Generative Models: GANs, Autoencoders and VAEs}
\companion{Computerphile: Generative Adversarial Networks (GANs)}{\yts{Computerphile+Generative+Adversarial+Networks+GANs}}

GANs are in the syllabus's architectures list. The InfoEdge hire had worked on GANs, so expect interviewers to push beyond the slogan
``a forger and a detective'' into the objective, the optimal discriminator, mode collapse and alternatives. Autoencoders and VAEs complete
the picture of how neural networks learn to \emph{generate}.

\section{Discriminative vs generative, once more}
A discriminative model learns $P(y\mid\mathbf x)$: ``is this profile photo real?''. A \textbf{generative model} learns the data distribution $P(\mathbf x)$
itself, so it can \emph{create} new samples: a realistic face, a property photo, a job description. Explicit models give a density (GMMs, VAEs
approximately, autoregressive LLMs); implicit models (GANs) only give a sampler.

\section{GANs: the forger and the detective}
\begin{itemize}
\item The \textbf{generator} $G$ takes random noise $\mathbf z$ (e.g. 100 numbers from a Gaussian) and outputs a fake sample $G(\mathbf z)$.
\item The \textbf{discriminator} $D$ takes a sample and outputs the probability it is real.
\item They are trained against each other: $D$ learns to tell real from fake; $G$ learns to fool $D$. As $D$ gets better, $G$ is forced to produce more
  realistic samples, until (ideally) fakes are indistinguishable.
\end{itemize}

\subsection{The minimax objective}
\[ \min_G\max_D\ V(D, G) = \E_{\mathbf x\sim p_{data}}[\log D(\mathbf x)] + \E_{\mathbf z\sim p_z}[\log(1 - D(G(\mathbf z)))]. \]
$D$ maximises: high $D(\mathbf x)$ on real data, low $D(G(\mathbf z))$ on fakes. This is exactly binary cross-entropy with real $=1$, fake $=0$.
$G$ minimises the second term: it wants $D(G(\mathbf z))$ high.

\subsection{The optimal discriminator}
For a fixed $G$ whose samples have density $p_g$, at each point $\mathbf x$ the discriminator maximises $p_{data}\log D + p_g\log(1 - D)$. Setting the derivative to zero:
\[ D^*(\mathbf x) = \frac{p_{data}(\mathbf x)}{p_{data}(\mathbf x) + p_g(\mathbf x)}. \]
\begin{example}[: optimal discriminator]
$p_{data}(\mathbf x) = 0.6$, $p_g(\mathbf x) = 0.2$: $D^* = 0.75$. At the global optimum $p_g = p_{data}$: $D^* = 0.5$ everywhere (the detective is reduced to guessing), and
$V(D^*, G^*) = \log0.5 + \log0.5 = -2\log2 = -1.386$.
\end{example}
Plugging $D^*$ in, $G$ minimises $-\log4 + 2\cdot\text{JSD}(p_{data}\,\|\,p_g)$: the \textbf{Jensen--Shannon divergence} between real and generated distributions
(maximum value $\log2 = 0.693$ when they do not overlap). So in the ideal case, training a GAN minimises JSD.

\subsection{Training loop}
Alternate: (1) update $D$ on a batch of real (label 1) and fake (label 0) samples; (2) update $G$ through $D$ (backprop through $D$ into $G$; $D$'s weights frozen in this
step). Discriminator loss for one real with $D = 0.9$ and one fake with $D = 0.2$: $-\ln0.9 - \ln0.8 = 0.105 + 0.223 = 0.328$.

\subsection{Non-saturating generator loss}
Early in training fakes are obviously fake: $D(G(\mathbf z))\approx0$, so $\log(1 - D(G(\mathbf z)))$ is flat and gives $G$ almost no gradient. In practice $G$ instead
\emph{maximises} $\log D(G(\mathbf z))$ (minimises $-\log D(G(\mathbf z))$): same fixed point, strong gradients when $G$ is losing. With $D(G(\mathbf z)) = 0.1$:
loss $= -\ln0.1 = 2.303$.

\section{Why GANs are hard to train}
\begin{itemize}
\item \textbf{Non-convergence}: two players with opposing objectives can oscillate rather than settle (a game, not a minimisation).
\item \textbf{Mode collapse}: $G$ produces only a few kinds of samples (all faces look alike) that currently fool $D$, ignoring the diversity of the data.
\item \textbf{Vanishing gradients}: a too-strong $D$ (accuracy 100\%) gives $G$ no useful signal; when real and fake distributions do not overlap, JSD is constant
  ($\log2$) and uninformative.
\item \textbf{Sensitivity} to architecture and hyperparameters; no single loss value tracks quality (both losses can look flat while samples improve).
\end{itemize}
\textbf{Remedies}: DCGAN guidelines (strided convolutions instead of pooling, BatchNorm, ReLU in $G$ with \textbf{tanh} output for pixels in $[-1, 1]$, LeakyReLU in $D$,
Adam with $\eta = 2\times10^{-4}$, $\beta_1 = 0.5$); \textbf{one-sided label smoothing} (real targets 0.9); \textbf{WGAN} (below); \textbf{spectral normalisation} of $D$
(controls its Lipschitz constant); gradient penalty; minibatch discrimination; two-timescale learning rates; progressive growing.

\subsection{Wasserstein GAN}
Replace JSD with the \textbf{Wasserstein (earth mover's) distance}: the minimum ``work'' to move probability mass from $p_g$ to $p_{data}$. Unlike JSD, it decreases smoothly
as distributions move closer even when they do not overlap, so $G$ always gets useful gradients. The discriminator becomes a \textbf{critic} outputting an unbounded score
(no sigmoid), constrained to be 1-Lipschitz (weight clipping in the original, gradient penalty in WGAN-GP); the critic is updated several times (e.g. 5) per generator
step. The critic's loss correlates with sample quality.

\section{Important variants}
\begin{itemize}
\item \textbf{DCGAN}: convolutional GAN; generator starts from a $4\times4$ map and doubles resolution with transposed convolutions (kernel 4, stride 2, padding 1):
  $4\to8\to16\to32\to64$.
\item \textbf{Conditional GAN}: feed a label or condition $y$ to both $G$ and $D$ to control what is generated (``generate a 2BHK living-room photo'').
\item \textbf{pix2pix}: paired image-to-image translation (sketch $\to$ photo) with an L1 loss plus adversarial loss.
\item \textbf{CycleGAN}: unpaired translation (day $\leftrightarrow$ night) using a \textbf{cycle-consistency} loss: translating there and back should recover the original.
\item \textbf{StyleGAN}: style-based generator, photorealistic faces; the source of many fake profile pictures.
\end{itemize}

\section{Evaluating generative models}
\begin{itemize}
\item \textbf{Inception Score (IS)}: high when each generated image is confidently classified (quality) and the classes are diverse (diversity). Ignores real data.
\item \textbf{Fr\'echet Inception Distance (FID)}: fit Gaussians to Inception features of real and generated images and compute the Fr\'echet distance between them;
  lower is better; captures quality and diversity; standard metric.
\item Precision/recall for generative models; human evaluation.
\end{itemize}

\section{Autoencoders}
An \textbf{autoencoder} compresses an input to a small code and reconstructs it: encoder $\mathbf z = f(\mathbf x)$, decoder $\hat{\mathbf x} = g(\mathbf z)$, trained to
minimise reconstruction error. The bottleneck forces it to keep the important structure. A linear autoencoder with squared error learns the PCA subspace; non-linear ones
learn curved manifolds.
Variants: \textbf{denoising} (reconstruct clean input from a corrupted one), \textbf{sparse}, \textbf{contractive}. Uses: compression, feature learning, \textbf{anomaly
detection} (normal data reconstruct well; anomalies have high reconstruction error), pre-training.
A plain autoencoder is a poor \emph{generator}: its code space has holes; random codes decode to garbage.

\section{Variational autoencoders (VAEs)}
\begin{itemize}
\item The encoder outputs a \emph{distribution} over codes, $q(\mathbf z\mid\mathbf x) = \mathcal N(\boldsymbol\mu(\mathbf x),\text{diag}\,\boldsymbol\sigma^2(\mathbf x))$.
\item Sample $\mathbf z = \boldsymbol\mu + \boldsymbol\sigma\odot\boldsymbol\epsilon$, $\boldsymbol\epsilon\sim\mathcal N(0, I)$: the \textbf{reparameterisation trick}, which moves
  the randomness into $\boldsymbol\epsilon$ so gradients flow through $\boldsymbol\mu$ and $\boldsymbol\sigma$.
\item Loss (negative ELBO) $=$ reconstruction error $+$ $KL(q(\mathbf z\mid\mathbf x)\,\|\,\mathcal N(0, I))$. The KL term pulls all codes toward a standard normal, filling the
  latent space without holes, so sampling $\mathbf z\sim\mathcal N(0, I)$ and decoding gives plausible new data.
\end{itemize}
VAEs train stably and give a latent space with meaningful interpolation, but samples tend to be blurrier than GANs'.

\section{GANs vs VAEs vs diffusion vs autoregressive}
\begin{center}\small
\begin{tabularx}{\linewidth}{l X X X}
\toprule
& Training & Samples & Notes\\
\midrule
GAN & adversarial, unstable & sharp, fast (one pass) & mode collapse; no likelihood\\
VAE & stable (ELBO) & blurrier & meaningful latent space; approximate likelihood\\
Diffusion & stable (denoising) & best quality and diversity & slow sampling (many steps); powers modern image generators\\
Autoregressive & stable (next-token) & excellent for text & sequential sampling; exact likelihood; LLMs\\
\bottomrule
\end{tabularx}
\end{center}
\textbf{Diffusion models} learn to reverse a gradual noising process: add Gaussian noise over $T$ steps until the image is pure noise; train a network to predict the noise at
each step; generate by starting from noise and denoising step by step.

\begin{practical}[: generative models at InfoEdge]
\begin{itemize}
\item \textbf{Defensive}: detecting GAN/diffusion-generated profile photos on Jeevansathi and fake recruiter profiles (artefact detectors, reverse image search, behaviour
  signals). Fraud is adversarial too: as detectors improve, generators adapt.
\item \textbf{Data augmentation}: synthetic tabular data (CTGAN) or images for rare classes; must be validated on real data.
\item \textbf{Anomaly detection}: autoencoder reconstruction error on recruiter behaviour.
\item \textbf{Text generation} (JD generator, content feed's ``content manifestation agent''): autoregressive LLMs, not GANs.
\end{itemize}
\end{practical}

\stuck{Google ML: GAN training}{https://developers.google.com/machine-learning/gan/training}
\section{Interview questions}

\begin{qa}{Explain how a GAN works, including the objective.}
\oneline{A generator maps random noise to samples and a discriminator classifies samples as real or fake; they play a minimax game where the discriminator maximises binary
cross-entropy on real vs fake and the generator tries to make the discriminator label its samples real; at the optimum the generator's distribution matches the data and the
discriminator outputs 0.5.}
\why Objective $\min_G\max_D\E[\log D(\mathbf x)] + \E[\log(1 - D(G(\mathbf z)))]$. Training alternates gradient steps. In practice $G$ uses the non-saturating loss
$-\log D(G(\mathbf z))$. With the optimal $D$, the generator minimises the Jensen--Shannon divergence.
\pushes \emph{``What does $G$ never see?''} Real data directly: it learns only through $D$'s gradients.
\end{qa}

\begin{qa}{Derive the optimal discriminator and compute it for $p_{data} = 0.6$, $p_g = 0.2$.}
\oneline{$D^* = p_{data}/(p_{data} + p_g) = 0.75$.}
\why Maximise $f(D) = a\log D + b\log(1 - D)$; $f'(D) = a/D - b/(1 - D) = 0\Rightarrow D = a/(a + b)$.
\end{qa}

\begin{qa}{What is the value of the GAN objective at the global optimum, and what divergence is being minimised?}
\oneline{$V = -\log4 = -1.386$, with $D^* = 0.5$ everywhere; the generator minimises the Jensen--Shannon divergence between data and generated distributions.}
\end{qa}

\begin{qa}{Why is the non-saturating generator loss used?}
\oneline{Because early on the discriminator easily rejects fakes, making $\log(1 - D(G(\mathbf z)))$ flat with tiny gradients; maximising $\log D(G(\mathbf z))$ has the same fixed point but
large gradients exactly when the generator is doing badly.}
\worked At $D(G(\mathbf z)) = 0.1$: gradient of $\log(1 - D)$ w.r.t.\ $D$ is $-1/0.9 = -1.11$; of $-\log D$ is $-10$.
\end{qa}

\begin{qa}{What is mode collapse and how do you mitigate it?}
\oneline{The generator produces only a narrow set of outputs that fool the current discriminator, ignoring other modes of the data; mitigations include WGAN/WGAN-GP, minibatch
discrimination, unrolled GANs, spectral normalisation, feature matching and diversity-promoting losses.}
\why The generator is rewarded for any sample that fools $D$, not for covering the data; $D$ then learns to reject that mode, $G$ jumps to another, and they cycle.
\end{qa}

\begin{qa}{Why are GANs unstable, and what did WGAN change?}
\oneline{They solve a two-player game rather than a minimisation, and with non-overlapping distributions the JS divergence gives no useful gradient; WGAN uses the Wasserstein distance with a
Lipschitz-constrained critic, which provides smooth, informative gradients and a loss that tracks quality.}
\end{qa}

\begin{qa}{\deep If the discriminator reaches 100\% accuracy early and stays there, what is wrong and what do you do?}
\oneline{The discriminator overpowers the generator, so the generator receives vanishing or uninformative gradients; weaken or regularise $D$ (fewer updates, label smoothing, noise on inputs,
spectral norm, dropout), use the non-saturating or Wasserstein loss, or strengthen $G$.}
\end{qa}

\begin{qa}{How are GANs evaluated?}
\oneline{With FID (distance between Gaussian fits of Inception features of real and generated images; lower is better), Inception Score (confident and diverse classifications), precision/recall
for generative models, and human judgement.}
\end{qa}

\begin{qa}{Conditional GAN, pix2pix and CycleGAN: what does each do?}
\oneline{A conditional GAN conditions both networks on extra information to control outputs; pix2pix does paired image-to-image translation with adversarial plus L1 loss; CycleGAN does unpaired
translation using cycle-consistency between two generators.}
\end{qa}

\begin{qa}{What is an autoencoder, and how can it detect anomalies?}
\oneline{A network that compresses an input to a low-dimensional code and reconstructs it; trained only on normal data, it reconstructs normal inputs well and anomalies poorly, so high reconstruction
error flags anomalies.}
\why Example: recruiter behaviour vectors; the top 0.1\% reconstruction errors go to the trust-and-safety queue. Choose the threshold on validation data with known incidents.
\end{qa}

\begin{qa}{\deep Why can a VAE generate new data but a plain autoencoder usually cannot?}
\oneline{A VAE's KL term forces codes to fill a known prior (a standard normal), so random draws from the prior decode to plausible samples; a plain autoencoder's codes occupy an irregular region
with gaps, so random codes decode to nonsense.}
\end{qa}

\begin{qa}{What is the reparameterisation trick?}
\oneline{Writing a random sample as a deterministic function of the parameters and independent noise, $\mathbf z = \boldsymbol\mu + \boldsymbol\sigma\odot\boldsymbol\epsilon$, so gradients can flow
through $\boldsymbol\mu$ and $\boldsymbol\sigma$ during backpropagation.}
\end{qa}

\begin{qa}{\deep GAN vs VAE vs diffusion: which would you pick to generate realistic synthetic property photos for augmentation?}
\oneline{A diffusion model (or fine-tuned pre-trained one) for the best quality and diversity, accepting slower sampling; a GAN if real-time sampling matters and you can manage training instability;
a VAE if you need a smooth latent space more than sharpness.}
\why For augmentation also ask whether synthetic images help at all: validate on real held-out photos; synthetic data can teach artefacts.
\end{qa}

\begin{qa}{Why does a DCGAN generator use tanh at the output?}
\oneline{Because images are scaled to $[-1, 1]$, and tanh produces outputs in exactly that range with zero-centred, non-saturating behaviour near the middle.}
\end{qa}

\begin{qa}{\deep How would you detect AI-generated fake profile photos on a matrimonial platform?}
\oneline{Combine an image classifier trained on real vs generated faces (with continual updates as generators change), artefact checks (eyes, backgrounds, frequency patterns), reverse image search
and duplicate detection, and behavioural signals (new account, many messages, copied bio text), routing high scores to human review.}
\why A single detector is brittle (adversarial adaptation); metadata and behaviour are harder to fake. Measure precision at the review queue's capacity; monitor drift.
\end{qa}

\begin{qa}{In one GAN training iteration, which network's weights change during the generator step?}
\oneline{Only the generator's; gradients flow through the discriminator to reach the generator, but the discriminator is not updated in that step.}
\end{qa}
